\section*{Chapter \ref{ch:b1:nucleophilic-substitution} --- Nucleophilic Substitution}

\begin{solution}{exo:b1:nucleophilic-substitution:1}
\ce{CH3CH2OH + Br-}; \ce{CH3OCH2CH3 + I-}; \ce{CH3CH2CH2CN + Cl-};
\ce{CH3NH3+ + Br-}. \omterm{def:b1:nucleophilic-substitution:substitution}{Substrates}: the halogenoalkanes; \omterm{def:b1:electronic-effects:nucleophile}{nucleophiles}:
\ce{OH-}, \ce{CH3CH2O-}, \ce{CN-}, \ce{NH3}; \omterm{def:b1:electronic-effects:nucleophile}{leaving groups}: the halide ions.
\end{solution}

\begin{solution}{exo:b1:nucleophilic-substitution:2}
Second order: tripled; both halved, divided by 4. First order in RY only:
unchanged; halved.
\end{solution}

\begin{solution}{exo:b1:nucleophilic-substitution:3}
SN2 (primary, good \omterm{def:b1:electronic-effects:nucleophile}{nucleophile}, aprotic); SN1 (tertiary, water); SN2
(methyl, strong \omterm{def:b1:electronic-effects:nucleophile}{nucleophile}); SN1 \omterm{def:b1:nucleophilic-substitution:sn1}{solvolysis} (secondary, weak neutral
\omterm{def:b1:electronic-effects:nucleophile}{nucleophile}, \omterm{def:b1:intermolecular-forces-solvents:solvent-classes}{protic solvent}), possibly with some SN2.
\end{solution}

\begin{solution}{exo:b1:nucleophilic-substitution:4}
The cyanide takes the place opposite the bromine and the three other
groups flip. In the product the CN group ranks first, as Br did, the
other ranks unchanged: the inverted arrangement has the opposite
descriptor, \cip{S}-2-methylbutanenitrile.
\end{solution}

\begin{solution}{exo:b1:nucleophilic-substitution:5}
1. \omterm{def:b1:electronic-effects:cleavage}{Heterolysis}: arrow from the C--Br bond to Br, giving \ce{(CH3)3C+} and
\ce{Br-} (slow). 2. A \omterm{def:b1:lewis-resonance-vsepr:lewis}{lone pair} of water to the cationic carbon:
\ce{(CH3)3C-OH2+}. 3. A \omterm{def:b1:lewis-resonance-vsepr:lewis}{lone pair} of another water molecule takes a proton
of the \ce{OH2+} group, the O--H pair returning to oxygen. Only the first,
slow step decides the rate; water, the solvent, is in such excess that
its concentration does not change.
\end{solution}

\begin{solution}{exo:b1:nucleophilic-substitution:6}
The carbon bearing Br is primary, but next to it a tert-butyl group fills
the space behind the C--Br bond: the SN2 \omterm{def:b1:elementary-steps:energy-profile}{transition state} is very
crowded. SN1 would need a primary \omterm{def:b1:electronic-effects:intermediates}{carbocation}, too unstable.
\end{solution}

\begin{solution}{exo:b1:nucleophilic-substitution:7}
Iodide replaces iodide by SN2, each time with inversion: \cip{S} becomes
\cip{R} until both are equally present, a \omterm{def:b1:stereochemistry-in-depth:enantiomers}{racemic mixture}. Each
substitution turns one \cip{S} molecule into one \cip{R}: the rotation
falls by two molecules' worth, so the rotation decays twice as fast as
the \omterm{def:b1:nucleophilic-substitution:substitution}{substrate} is substituted.
\end{solution}

\begin{solution}{exo:b1:nucleophilic-substitution:8}
SN2: bromomethane $>$ 1-bromopropane $>$ 2-bromopropane $>$
2-bromo-2-methylpropane (crowding). SN1: the reverse order (\omterm{def:b1:electronic-effects:intermediates}{carbocation}
stability), bromomethane and 1-bromopropane hardly reacting.
\end{solution}

\begin{solution}{exo:b1:nucleophilic-substitution:9}
$x_{\text{inv}} - x_{\text{ret}} = 0.12$ and $x_{\text{inv}} + x_{\text{ret}} = 1$:
\qty{56}{\percent} inverted, \qty{44}{\percent} retained. Mostly SN1
\omterm{def:b1:nucleophilic-substitution:sn1}{racemisation}, with a slight excess of inversion: the departing bromide
still shields its side when water arrives.
\end{solution}

\begin{solution}{exo:b1:nucleophilic-substitution:10}
$v = k_1k_2[\mathrm{RY}][\mathrm{Nu}]/(k_{-1}[\mathrm{Y^-}] + k_2[\mathrm{Nu}])$
(\cref{prop:b1:nucleophilic-substitution:sn1-law}). Adding \ce{Y-}
increases the denominator: more \omterm{def:b1:electronic-effects:intermediates}{carbocations} return to the \omterm{def:b1:nucleophilic-substitution:substitution}{substrate}
instead of being captured, and the rate falls. Negligible when
$k_2[\mathrm{Nu}] \gg k_{-1}[\mathrm{Y^-}]$, as when the \omterm{def:b1:electronic-effects:nucleophile}{nucleophile} is the
solvent.
\end{solution}

\begin{solution}{exo:b1:nucleophilic-substitution:11}
Equal concentrations: $-\mathrm{d}c/\mathrm{d}t = kc^2$, $1/c = 1/C_0 + kt$,
$t_{1/2} = 1/(kC_0) = \qty{5.0e5}{s}$ (about six days). With a fifty-fold
excess of hydroxide, $[\ce{OH-}] \approx 50C_0$ is constant: $c = C_0
e^{-k't}$ with $k' = 50kC_0 = \qty{1.0e-4}{s^{-1}}$, $t_{1/2} = \ln 2/k' =
\qty{6.9e3}{s}$, about two hours.
\end{solution}

\begin{solution}{exo:b1:nucleophilic-substitution:12}
Bromide would have to expel \ce{OH-}, a \omterm{def:b1:predominant-reaction:strong-acid}{strong base} and a poor leaving
group: no reaction. In concentrated HBr the OH is protonated to
\ce{-OH2+}, whose \omterm{def:b1:electronic-effects:nucleophile}{leaving group} is water. Tertiary alcohol: water leaves,
giving the \omterm{def:b1:electronic-effects:intermediates}{carbocation}, captured by \ce{Br-} (SN1). Primary alcohol:
\ce{Br-} attacks the carbon of \ce{R-OH2+} from behind as water departs
(SN2).
\end{solution}

\begin{solution}{pb:b1:nucleophilic-substitution:1}
\textbf{1.} Each molecule that reacts gives one \ce{H+} and one \ce{Cl-};
the conductivity is the sum of the ionic contributions, proportional to
the amount of ions.
\textbf{2.} $[\mathrm{RCl}]/[\mathrm{RCl}]_0 = (\kappa_\infty -
\kappa)/\kappa_\infty$.
\textbf{3.} $\ln[\mathrm{RCl}] = \ln[\mathrm{RCl}]_0 - kt$: plot
$\ln(\kappa_\infty - \kappa)$ against $t$.
\textbf{4.} 0.693, 0.603, 0.513, 0.423, 0.333, 0.153, $-0.117$, $-0.387$.
\textbf{5.} Yes, slope $-0.0180$ per minute throughout: first order.
\textbf{6.} Any order in water: its concentration is constant
(degeneracy of the order).
\textbf{7.} $k = 0.0180/60 = \qty{3.0e-4}{s^{-1}}$.
\textbf{8.} $t_{1/2} = \ln 2/k = \qty{2.3e3}{s}$, \qty{38.5}{min}.
\textbf{9.} $\ln 10/k = \qty{7.7e3}{s}$, \qty{128}{min}.
\textbf{10.} $\ln(\kappa_\infty - \kappa)$: 0.693 at $t = 0$, 0.109 at
\qty{10}{min}, and so on: slope $-0.0584$ per minute, $k = \qty{9.7e-4}{s^{-1}}$.
\textbf{11.} $9.7/3.0 = 3.2$.
\textbf{12.} At \qty{30}{min}: $2.000(1 - e^{-0.0584 \times 30}) = 1.653$,
measured 1.654.
\textbf{13.} SN1: tertiary \omterm{def:b1:nucleophilic-substitution:substitution}{substrate}, \omterm{def:b1:intermolecular-forces-solvents:solvent-classes}{protic solvent}, neutral \omterm{def:b1:electronic-effects:nucleophile}{nucleophile}.
\textbf{14.} As in \cref{exo:b1:nucleophilic-substitution:5}, with Cl.
\textbf{15.} The slow ionisation involves only the \omterm{def:b1:nucleophilic-substitution:substitution}{substrate}: $v =
k[\mathrm{RCl}]$.
\textbf{16.} Hydroxide could only act in the fast capture step, after the
rate-determining ionisation; the rate does not depend on it.
\textbf{17.} The \omterm{def:b1:electronic-effects:intermediates}{carbocation} is planar and \omterm{def:b1:stereochemistry-in-depth:stereocentre}{achiral}; water adds on both
faces equally: racemic alcohol.
\textbf{18.} A high first barrier (ionisation) to the \omterm{def:b1:electronic-effects:intermediates}{carbocation} in a
shallow well, then a low barrier to capture: the first step is
rate-determining.
\textbf{19.} It would form a primary \omterm{def:b1:electronic-effects:intermediates}{carbocation} (too unstable), and water
is too weak a \omterm{def:b1:electronic-effects:nucleophile}{nucleophile} for a fast SN2.
\textbf{20.} $k = A\,e^{-E_a/RT}$.
\textbf{21.} $E_a = R\ln(k_2/k_1)/(1/T_1 - 1/T_2)$.
\textbf{22.} $E_a = 8.314 \times \ln(9.74/3.00)/(1/298.15 - 1/308.15) =
\qty{9.0e4}{J/mol}$.
\textbf{23.} The height of the barrier of the ionisation, the
\omterm{def:b1:elementary-steps:rds}{rate-determining step}.
\textbf{24.} $k = 3.0 \times 10^{-4} \exp\bigl(\frac{90\,000}{8.314}
(\frac{1}{298.15} - \frac{1}{318.15})\bigr) = \qty{2.9e-3}{s^{-1}}$.
\textbf{25.} \textbf{$E_a = \qty{90}{kJ/mol}$}.
\end{solution}
